% GTDynamics Copyright 2026, Georgia Tech Research Corporation.
% See LICENSE for license information.
% Build from the repository root:
%   mkdir -p build/doc/cmc-admm
%   latexmk -pdf -outdir=build/doc/cmc-admm doc/cmc-admm-linear-step.tex
%   cp build/doc/cmc-admm/cmc-admm-linear-step.pdf doc/
\documentclass[11pt,letterpaper]{article}
\usepackage[margin=0.9in]{geometry}
\usepackage[T1]{fontenc}
\usepackage{lmodern}
\usepackage{amsmath,amssymb,mathtools}
\usepackage{booktabs,tabularx}
\usepackage{microtype}
\usepackage{xcolor}
\usepackage{enumitem}
\usepackage{fancyhdr}
\usepackage[colorlinks=true,linkcolor=blue!45!black,urlcolor=blue!45!black,
  citecolor=blue!45!black]{hyperref}
\hypersetup{pdftitle={The ADMM Linear Step in CMC-Opt},
  pdfauthor={GTDynamics contributors}}
\newcommand{\R}{\mathbb{R}}
\newcommand{\diag}{\operatorname{diag}}
\newcommand{\pos}[1]{[#1]_{+}}
\newcommand{\code}[1]{\texttt{\small #1}}
\setlength{\parindent}{0pt}
\setlength{\parskip}{5pt plus 1pt}
\setlist{nosep,leftmargin=*}
\pagestyle{fancy}
\fancyhf{}
\fancyhead[L]{\small GTDynamics technical note}
\fancyhead[R]{\small CMC-Opt / ADMM}
\fancyfoot[C]{\thepage}
\setlength{\headheight}{14pt}
\title{\vspace{-1.7em}The ADMM Linear Step in CMC-Opt}
\author{GTDynamics contributors}
\date{November 12, 2026}

\begin{document}
\maketitle
\thispagestyle{fancy}
\vspace{-1.6em}
\begin{abstract}
CMC-Opt reduces equality-constrained robot states to local manifold coordinates
and solves a global tangent-cone quadratic program at each Levenberg--Marquardt
trial. This note describes an experimental ADMM backend for that linear step.
Constraint-output splitting and diagonal scaling turn the inequality update
into scalar clipping, while GTSAM's multifrontal solver supplies a factorization
reused across all ADMM iterations of a fixed QP. Acceptance uses an optimality
certificate in the original QP coordinates; nonlinear retraction remains a
separate CMC-Opt operation.
\end{abstract}

\section{From CMC-Opt to a reduced quadratic program}
CMC-Opt optimizes a factor-graph objective over constraint manifolds with
inequality boundaries and corners. Its nonlinear problem has the form
\begin{equation}
  \min_x \frac12\sum_a\|f_a(x_a)\|_{\Sigma_a}^{2}
  \quad\text{subject to}\quad h_i(x_i)=0,\quad c_i(x_i)\leq0,
  \label{eq:nlp}
\end{equation}
where $\|v\|_\Sigma^2=v^\top\Sigma^{-1}v$. Constraint-connected variables form
local components, but objective factors can couple different components and
times. The detailed geometry and retraction machinery are documented in the
repository's CMC-Opt math-to-code map and component-coordinate guide
\cite{cmc,coordinates}.

Equality elimination happens before ADMM sees the problem. At the current
feasible state, component $i$ provides a basis $B_i$ satisfying
$\mathrm{D}h_i(x_i)B_i=0$. Write its ambient tangent increment as
$\delta_i=B_i\xi_i$; unconstrained keys retain their ordinary tangent
coordinates. With $B$ denoting the assembled local lifting maps, the linear
step is
\begin{equation}
  \min_{\xi\in\R^n} q_\lambda(\xi)
  =\frac12\|JB\xi+r\|_W^2+\frac\lambda2\|L\xi\|^2
  =\frac12\xi^\top H\xi+g^\top\xi+c,
  \qquad G\xi\leq0.
  \label{eq:qp}
\end{equation}
Here $\|v\|_W^2=v^\top Wv$ uses the cost information weights, $L$ is the
chosen LM damping operator, and $H$ includes the current damping. The QP has
\emph{no explicit equality constraints}. Each block $G_i$ contains the
appropriately weighted derivatives of selected boundary inequalities,
$\mathrm{D}c_{i,\mathcal A_i}(x_i)B_i$, and touches only $\xi_i$.

The homogeneous inequalities encode the current tangent cone, rather than all
affine linearizations of the original bounds. Slack inequalities are omitted;
selected near-boundary modes use the same homogeneous construction. Nonlinear
retraction handles boundaries encountered by a finite step. Thus zero is
feasible for the QP, but QP feasibility alone does not establish feasibility
of the final nonlinear state.

\newpage
\section{Choosing the split and the scaling}
\subsection{Manifold-variable splitting}
The initial ADMM formulation duplicates each reduced component variable.
Define $C_i=\{v:G_i v\leq0\}$ and enforce $\xi_i=z_i\in C_i$.
For a positive diagonal metric $M$ and scaled dual $u$, its updates are
\begin{align}
  (H+\rho M)\xi^{k+1}&=-g+\rho M(z^k-u^k),\label{eq:point}\\
  z_i^{k+1}&=\mathop{\arg\min}_{v\in C_i}
       \tfrac12\|v-(\xi_i^{k+1}+u_i^k)\|_{M_i}^2,\nonumber\\
  u^{k+1}&=u^k+\xi^{k+1}-z^{k+1}.\nonumber
\end{align}
The implementation uses local dual-coordinate/Dykstra projections. The
identity metric is the default; an optional metric uses the floored diagonal
of $H$. This is already a split by CMC components, not one independent
optimization problem per scalar coordinate. The globally coupled solve
remains necessary because objective factors link the components.

\subsection{Constraint-output splitting}
The newer formulation assigns one auxiliary scalar to each tangent inequality
row, so the local update becomes an orthant projection. Introduce a positive
diagonal variable scaling $D$ and row scaling $E$, and define
\begin{equation}
  \xi=Dy,\qquad \bar H=DHD,\qquad \bar g=Dg,\qquad \bar G=EGD.
  \label{eq:scaled}
\end{equation}
The implemented scaling, with $\varepsilon=10^{-12}$, is
\begin{equation}
  D_{jj}=\frac{1}{\sqrt{\max(H_{jj},\varepsilon)}},\qquad
  E_{\ell\ell}=\begin{cases}
    1/\max(\|(GD)_{\ell,:}\|_2,\varepsilon),
      &\|(GD)_{\ell,:}\|_2>0,\\
    1,&\text{otherwise}.
  \end{cases}
  \label{eq:scales}
\end{equation}
Both scalings are fixed within a QP and recomputed when damping changes.
With scaling disabled, $D=I$ and $E=I$. The diagonal floor defines the
coordinate transformation; it does not add a regularizer to $H$.

The equivalent split problem is
\begin{equation}
  \min_{y,z}\ \tfrac12 y^\top\bar H y+\bar g^\top y
       +\iota_{\R_-^m}(z)
  \quad\text{subject to}\quad \bar G y-z=0,
  \label{eq:split}
\end{equation}
where $m$ is the number of selected scalar inequality rows and the indicator
is zero on the nonpositive orthant and infinite elsewhere. Positive $E$ and
invertible $D$ preserve the original feasible set and objective under the
coordinate change. They change the ADMM iteration geometry and hence the
meaning of a given numerical $\rho$.

Scaling matters because reduced coordinates need not be orthonormal or
dimensionless. An elimination basis can retain physical joint, acceleration,
and force coordinates with very different curvature. Diagonal equilibration
addresses that imbalance, but does not remove all correlations or guarantee
a well-conditioned system. The broader sensitivity of ADMM to scaling is
discussed in~\cite{scaling}; the specific rule above is the implemented heuristic.

\newpage
\section{The ADMM iteration and the cached GTSAM solve}
Scaled ADMM applies to~\eqref{eq:split} with fixed $\rho>0$ and the scaled
dual $u$~\cite{admm}. Completing the square in its augmented Lagrangian gives
\begin{equation}
  \mathcal L_\rho(y,z,u)=\tfrac12y^\top\bar H y+\bar g^\top y
   +\iota_{\R_-^m}(z)+\tfrac\rho2\|\bar Gy-z+u\|^2
   -\tfrac\rho2\|u\|^2.
\end{equation}
Each iteration consists of a global linear solve, scalar clipping, and a
dual update:
\begin{align}
  K y^{k+1}&=-\bar g+\rho\bar G^\top(z^k-u^k),
       &K&=\bar H+\rho\bar G^\top\bar G,\label{eq:global}\\
  z^{k+1}&=\min(\bar G y^{k+1}+u^k,0),&&\label{eq:clip}\\
  u^{k+1}&=u^k+\bar G y^{k+1}-z^{k+1}.&&\label{eq:dual}
\end{align}
The minimum is componentwise. The current solver initializes $z^0=u^0=0$;
the first global solve therefore has right-hand side $-\bar g$. A QP with
no inequality rows reduces to one direct solve. The intended input is a
positive-definite damped objective; the solver adds no hidden regularization.

\subsection{Retaining the factor graph}
GTSAM factors the augmented Gaussian graph without assembling a global dense
Hessian or transferring the system to an external sparse solver. Each original
reduced Jacobian factor receives the column transformation $A_i\mapsto A_iD_i$,
with its right-hand side and noise model preserved. Local Hessian-factor
inputs retain their Hessian representation under the same change of
coordinates. For each inequality block the graph gains the unary factor
\begin{equation}
  \sqrt\rho\,\bar G_i y_i\approx0.
\end{equation}
These penalty factors change diagonal variable blocks but introduce no
new edges between CMC components. Original objective couplings still
determine the global graph and its elimination fill.

The implementation calls \code{MultifrontalSolver::eliminateInPlace} and
\code{computeBayesTree} once per QP. Abstractly, the resulting Bayes tree
represents $K=R^\top R$ in the elimination ordering. For a new right-hand side
$b^k$ from~\eqref{eq:global}, the update uses
\begin{equation}
  R^\top v=b^k,\qquad R y^{k+1}=v.
\end{equation}
These operations are \code{backSubstituteTranspose} followed by
\code{backSubstitute}; the cached numerical factors are unchanged.

\subsection{What can be reused}
Every ADMM iteration of a fixed QP reuses the numerical factorization because
$H,G,D,E,\rho$ are fixed. A change only to the objective's linear term or to
$z-u$ changes the right-hand side, so triangular solves suffice. Changing LM
damping changes $H$, generally $D$ and $E$, and hence $K$: it requires a
\emph{new numerical factorization}. The workspace reuses symbolic analysis,
ordering, and clique storage across structurally compatible damping trials.
A new nonlinear linearization creates a new workspace. Changing $\rho$
would also require refreshing the numerical factors.

\newpage
\section{Certification and return to the nonlinear optimizer}
ADMM convergence is certified against the original reduced QP, rather than
only the scaled consensus equation. Transform the candidate and inequality
multipliers back as
\begin{equation}
  \xi=Dy,\qquad \mu=\rho E u.
  \label{eq:recover}
\end{equation}
The multiplier transformation follows by equating the linear constraint terms
$\mu^\top G\xi=(\rho u)^\top EG\xi$. It is essential when row scaling is
enabled. The original-coordinate stationarity residual is
\begin{equation}
  s=\nabla q_\lambda(\xi)+G^\top\mu
    =H\xi+g+G^\top\mu.
  \label{eq:stationarity}
\end{equation}
The implementation evaluates the gradient from the original Gaussian factors,
so the certificate includes the original weights and the current LM damping.
``Original coordinates'' here means the unscaled, equality-reduced QP, not
the redundant ambient coordinates eliminated by $B$.

\subsection{Original-coordinate KKT checks}
The solver accepts only finite candidates and finite objective values satisfying
the following tests, with $\tau=10^{-7}$:
\begin{align}
  \|\pos{G\xi}\|_\infty&\leq\tau
      &&\text{(primal feasibility)},\\
  \mu_\ell&\geq-\tau\quad\text{for every row }\ell
      &&\text{(dual feasibility)},\\
  \|\mu\odot(G\xi)\|_\infty&\leq\tau
      &&\text{(complementarity)},\\
  \|s\|_\infty&\leq\tau
      &&\text{(stationarity)}.
\end{align}
Checks occur at iterations $1,26,51,\ldots$ and at the iteration cap. For the
convex QP these are numerical KKT tests for optimality; their absolute
tolerance is an implementation choice, not a unit-invariant error bound.
Failure to certify within the limit rejects the LM trial.

The recorded primal residual is $\|\bar Gy-z\|_\infty$. In constraint-split
mode the field named \code{dualResidual} records the latest original-coordinate
stationarity norm, rather than the textbook ADMM residual
$\rho\bar G^\top(z^k-z^{k-1})$. This mode has no additional consensus-residual
acceptance gate beyond the KKT certificate above.

\subsection{Active faces and retraction}
The solver returns a blocking-row mask in the original QP row order. Row
$\ell$ is marked active when $\mu_\ell>0$ and
$|(G\xi)_\ell|\leq\tau$. A tight row with zero multiplier need not be marked.
The CMC adapter promotes a nonlinear vector constraint block to a retraction
equality only when all of its rows are active; a partially active block
remains an inequality.

The outer optimizer lifts the step through $B$ and performs its existing
directional checks and nonlinear retraction. Retraction restores the nonlinear
equalities and inequalities, including omitted slack boundaries that the
finite step reaches. The undamped model determines predicted reduction;
nonlinear objective evaluation and LM acceptance decide whether to keep the
trial. ADMM solves only~\eqref{eq:qp}, not the nonlinear feasibility problem
or the retraction subproblem.

\newpage
\section{Cost, implementation, and scope}
The solve time includes both factorization and repeated Bayes-tree solves:
\begin{equation}
  T_{\mathrm{QP}}\approx T_{\mathrm{setup+factor}}
    +N\bigl(T_{\mathrm{two\ triangular\ solves}}+T_{G,G^\top}+O(m)\bigr)
    +N_{\mathrm{check}}T_{\mathrm{KKT}}.
\end{equation}
Triangular-solve cost depends on clique dimensions and fill; the iteration
count $N$ depends on curvature, constraint geometry, scaling, $\rho$, and
accuracy. Constraint splitting removes iterative cone projections, but
retains the global primal solve at each iteration.

In the G1 timing example, 171 CMC blocks contain 18,297 scalar coordinates;
42 unconstrained six-dimensional blocks add 252. The total of 18,549 is
\emph{after equality elimination}. The captured second QP has 230 inequality
rows on 117 CMC blocks: $z$ and $u$ each have dimension 230, while the global
solve still has dimension 18,549. These counts are specific to that capture.

\subsection{Code and options}
The constraint-output path shares the historical point-split backend API:
\begin{center}
\small
\begin{tabularx}{\linewidth}{@{}p{0.47\linewidth}X@{}}
\toprule
File under \code{gtdynamics/cmcopt/} & Responsibility \\
\midrule
\code{ConstraintSplitAdmmSolver.cpp} & Scaling, augmented graph, cached solves,
clipping, multiplier recovery, and KKT certificate. \\
\code{PointSplitAdmmSolver.h/.cpp} & Shared workspace and validation;
original component-variable splitting. \\
\code{IELMOptimizer.cpp} & Workspace lifetime across damping trials. \\
\code{IELMOptimizerState.cpp} & Damped QP assembly, backend dispatch,
blocking-row mapping, and nonlinear update. \\
\bottomrule
\end{tabularx}
\end{center}
Select \code{QpBackend::PointSplitADMM} and choose \code{admmRho}.\\
Enable \code{admmConstraintSplitting} and \code{admmUseDiagonalMetric};
the latter flag selects $D,E$ scaling here and the $M$-metric in point mode.
This opt-in experiment uses fixed $\rho$, cold starts for each QP, and no
over-relaxation, polishing, or fallback. Tests are in \code{testPointSplitAdmm}
and \code{testIELMQp}. Compare timings only with comparable full QP convergence
and feasible nonlinear results~\cite{experiment}.

\begingroup
\small
\raggedright
\setlength{\parskip}{0pt}
\begin{thebibliography}{5}
\bibitem{cmc} GTDynamics contributors.
\emph{CMC-Opt Math to Code Map}. Repository document,
\path{gtdynamics/cmcopt/README.md}; includes pointers to Chapter 4 of
Yetong Zhang's 2024 thesis on constraint manifold optimization.
\bibitem{coordinates} GTDynamics contributors.
\emph{Equality coordinates and the global tangent-cone QP}.
\path{gtdynamics/cmcopt/ComponentCoordinates.md}.
\bibitem{admm} S. Boyd, N. Parikh, E. Chu, B. Peleato, and J. Eckstein.
\href{https://web.stanford.edu/~boyd/papers/admm_distr_stats.html}
{Distributed Optimization and Statistical Learning via the Alternating
Direction Method of Multipliers}.
\emph{Foundations and Trends in Machine Learning}, 3(1):1--122, 2011.
\bibitem{scaling} P. Giselsson and S. Boyd.
\href{https://web.stanford.edu/~boyd/papers/diag_scaling_DR_ADMM.html}
{Diagonal Scaling in Douglas--Rachford Splitting and ADMM}.
\emph{IEEE Conference on Decision and Control}, pp.\ 5033--5039, 2014.
\bibitem{experiment} GTDynamics contributors.
\emph{Point-splitting ADMM versus ProxQP (experimental)}.
\path{examples/cmc_benchmarks/PointSplitADMM.md}.
\end{thebibliography}
\endgroup
\end{document}
